\section{Experimental Protocol}
\label{sec:method}

\paragraph{Models and decoding.}
We evaluate Qwen2.5-7B-Instruct\footnote{\url{https://huggingface.co/Qwen/Qwen2.5-7B-Instruct}} and Llama-3.1-8B-Instruct\footnote{\url{https://huggingface.co/meta-llama/Llama-3.1-8B-Instruct}}, both 4-bit quantised, with greedy decoding.
The model reasons first and ends with its verdict and stated confidence.
The order matters: in a pilot with the verdict first, the models often contradicted their own working on MATH, so the verdict was close to a first guess.
Qwen generates up to 768 new tokens per response and Llama up to 1{,}024.

\paragraph{Verdict recovery.}
Llama often fell into repetition loops (``However, I must correct my previous response \dots'') that ran into the token limit without a verdict; in a MATH run without a fallback \LlamaMathOldMissing\% of its turns had none, rising with the turn as looping answers filled the context.
For Llama we therefore apply a fallback: if a response has no parseable verdict, a sentence repeated three or more times is cut back to one copy, then ``\texttt{[Verdict:}'' is appended and the model completes the line, a two-stage answer extraction in the style of \citet{kojima2022large}.
Later turns see the completed answer.
On the hardest MATH items, verdicts the model wrote itself agreed with the run without the fallback on all 22 turns where both had one.
Forced verdicts make up \LlamaMathForcedCtrl--\LlamaMathForcedCrit\% of Llama's MATH turns and \LlamaRAForcedCrit--\LlamaRAForcedCtrl\% of its real-analysis turns, so every Llama result is also reported with them treated as missing.
Qwen ran without the fallback; \QwenMathMissing\% of its MATH turns and \QwenRAMissing\% of its real-analysis turns have no verdict and are treated as missing.

\paragraph{Conversation.}
Turn 0 poses the neutral question; turns 1--6 send the six prompts of one track in an order shuffled per item with a fixed seed, keeping the whole history in context (Figure~\ref{fig:overview}b).
Every item is followed for all six turns on all three tracks.

\paragraph{Measurements.}
After each turn we record:
\begin{enumerate}
\setlength\itemsep{0pt}
\item the \emph{verdict}, parsed from the last verdict line;
\item the \emph{pre-answer probe}: a verdict slot (``\texttt{Conclusion: The statement is strictly [Verdict: }'') is appended to the conversation before the turn-$t$ response, and the next-token probabilities of ``Valid'' and ``Invalid'' are normalised over the two;
\item the \emph{self-check}: after the response, a further user turn asks whether the statement is (A) \Valid or (B) \Invalid, and we compare the next-token probabilities of \texttt{A} and \texttt{B};
\item the \emph{stated confidence}, the verbalised percentage \citep{tian-etal-2023-just, xiong2023express};
\item on MATH, the \emph{derived answer}: the integer the model's working reaches (the last boxed or stated result), with sentences that restate the claim removed.
\end{enumerate}
Each probe is read as favouring \Valid when its probability exceeds 0.5.

\paragraph{Tone severity.}
The nominal valences in Table~\ref{tab:categories} are our own.
We score each of the \ToneWrappers{} feedback wordings, with the claim removed, using established resources rather than a learned scorer (Appendix~\ref{app:scorer}):
mean valence and arousal of its words from the NRC VAD lexicon \citep{mohammad-2018-obtaining};
the probability that it is polite, from a logistic regression on ConvoKit politeness strategies trained on the Wikipedia requests of the Stanford Politeness Corpus \citep{danescu-niculescu-mizil-etal-2013-computational};
and the probabilities of negative (disapproval, annoyance, anger, disappointment) or positive (admiration, approval, pride, optimism) emotions from a classifier fine-tuned on GoEmotions \citep{demszky-etal-2020-goemotions}.
Within each track we orient the measures so that higher means stronger, z-score them, and average them into one \emph{severity} score (politeness is left out for praise, since flattery and a polite nudge are both polite).
The lexicon's dominance score was dropped after a first look at the scores: it rates the words' referents (``garbage'' is a low-dominance word), not the speaker's force, and ran against the categories ($\rho = \ToneDomRho$).
For each feedback turn whose previous verdict differs from what the feedback pushes for, we regress whether the model switches (logistic) on severity, with turn and difficulty as categorical controls and length as a covariate.

\paragraph{Statistics.}
The \emph{pressure effect} is the difference between a feedback track and the control on the same items at turn 6.
A claim and its twin share a problem, so all 95\% intervals are percentile bootstrap intervals that resample claim--twin \emph{pairs} (2{,}000 resamples for the main effects, 1{,}000 for the tone regressions), paired across tracks.
